%%%%%%%%%%%%%%%%%%%%%%%%%%%%%%%%%%%%%%%%%%%%%%%%%%%%%%%%%%%%%%%%%%%%%%%%%%%%%%%%
%%  content/05-evaluation.tex — فصل ۵: ارزیابی و نتایج
%%
%%  نمودار ستونی این فصل با TikZ خالص کشیده شده است تا به pgfplots نیازی نباشد.
%%  ترفند مقیاس: با y=0.032cm هر واحد محور عمودی برابر یک درصد است، پس مقدارها را
%%  می‌توان مستقیم به‌صورت درصد نوشت و ۱۰۰ درصد ۳٫۲ سانتی‌متر بلند می‌شود.
%%%%%%%%%%%%%%%%%%%%%%%%%%%%%%%%%%%%%%%%%%%%%%%%%%%%%%%%%%%%%%%%%%%%%%%%%%%%%%%%
\section{ارزیابی و نتایج}

%%%%%%%%%%%%%%%%%%%%%%%%%%%%%%%%%%%%%%%%%%%%%%%%%%%%%%%%%%%%%%%%%%%%%%%%%%%%%%%%
\begin{frame}{روش ارزیابی}
  \begin{description}
    \item[مجموعهٔ آزمون] \hl{۴۰ پرسش} در چهار دستهٔ رخداد تازه، پرسش چندبخشی،
          پرسش عددی و پرسش تعریفی.
    \item[سنجه‌ها] درستی پاسخ (داوری انسانی)، \en{Precision@3} و \en{NDCG@3}
          برای مرحلهٔ بازیابی.
    \item[مبنای مقایسه] یک سامانهٔ تجاری پرسش و پاسخ مبتنی بر وب، با همان
          پرسش‌ها و در همان بازهٔ زمانی.
    \item[مطالعهٔ کاهشی] اجرای زنجیره یک‌بار \term{بدون} بازرتبه‌بندی و یک‌بار
          \term{با} آن، برای جدا کردن سهم این مرحله.
  \end{description}

  \begin{PTnote}
    داوری روی درستی پاسخ به‌صورت دوسویه و بدون آگاهی از نام سامانه انجام شد تا
    سوگیری داور کنترل شود.
  \end{PTnote}
\end{frame}

%%%%%%%%%%%%%%%%%%%%%%%%%%%%%%%%%%%%%%%%%%%%%%%%%%%%%%%%%%%%%%%%%%%%%%%%%%%%%%%%
\begin{frame}{مقایسه با سامانهٔ مرجع}{۴۰ پرسش، داوری انسانی}
  \begin{center}
    \small
    \begin{tabular}{@{}lcc@{}}
      \toprule
      سنجه & سامانهٔ پیشنهادی & سامانهٔ مرجع \\
      \midrule
      درستی پاسخ            & \hl{۹۵٪}   & ۷۷٫۵٪ \\
      پاسخ همراه ارجاع      & ۱۰۰٪       & ۴۵٪   \\
      میانگین اسناد بافتار  & \hl{۳}     & ۱۰ تا ۲۰ \\
      \bottomrule
    \end{tabular}
  \end{center}

  \begin{PTresult}
    سامانهٔ پیشنهادی با بافتاری \hl{کوتاه‌تر}، درستی بالاتری به دست می‌آورد؛ پس
    بهبود از \term{کیفیت بازیابی} می‌آید، نه از حجم بافتار.
  \end{PTresult}
\end{frame}

%%%%%%%%%%%%%%%%%%%%%%%%%%%%%%%%%%%%%%%%%%%%%%%%%%%%%%%%%%%%%%%%%%%%%%%%%%%%%%%%
\begin{frame}{اثر بازرتبه‌بندی بر کیفیت بازیابی}
  \begin{center}
    \begin{tikzpicture}[
        x = 1cm, y = 0.032cm,
        PTbarA/.style = {fill=PTaccentdark, draw=PTaccentdark, line width=.4pt},
        PTbarB/.style = {fill=PTline,       draw=PTmuted,      line width=.4pt},
        PTgrid/.style = {draw=PTline, line width=.4pt,
                         dash pattern=on 1pt off 1.4pt},
        PTval/.style  = {font=\tiny, text=PTink,  anchor=south, inner sep=1pt},
        PTswatch/.style = {inner sep=0pt, minimum width=0.34cm,
                           minimum height=0.20cm},
      ]
      % ---- شبکه و محورها ------------------------------------------------------
      \draw[PTgrid] (1.6, 25) -- (11.4, 25);
      \draw[PTgrid] (1.6, 50) -- (11.4, 50);
      \draw[PTgrid] (1.6, 75) -- (11.4, 75);
      \draw[PTgrid] (1.6,100) -- (11.4,100);
      \draw[draw=PTmuted,line width=.6pt] (11.4,0) -- (11.4,106);
      \draw[draw=PTmuted,line width=.6pt] ( 1.6,0) -- (11.4,  0);

      % ---- برچسب محور عمودی: سمت راست، چون خواندن از راست آغاز می‌شود ---------
      \node[PTcapt,anchor=west] at (11.55,  0) {\PTrl{۰}};
      \node[PTcapt,anchor=west] at (11.55, 25) {\PTrl{۲۵}};
      \node[PTcapt,anchor=west] at (11.55, 50) {\PTrl{۵۰}};
      \node[PTcapt,anchor=west] at (11.55, 75) {\PTrl{۷۵}};
      \node[PTcapt,anchor=west] at (11.55,100) {\PTrl{۱۰۰}};

      % ---- گروه ۱ (راست): Precision@3 ---------------------------------------
      \fill[PTbarB] (7.48,0) rectangle (8.53,67.7);
      \fill[PTbarA] (8.67,0) rectangle (9.72,93.3);
      \node[PTval] at (8.005,67.7) {\PTrl{۶۷٫۷٪}};
      \node[PTval] at (9.195,93.3) {\PTrl{۹۳٫۳٪}};
      \node[PTcapt,anchor=north] at (8.6,-4) {\en{Precision@3}};

      % ---- گروه ۲ (چپ): NDCG@3 ----------------------------------------------
      \fill[PTbarB] (3.28,0) rectangle (4.33,82.4);
      \fill[PTbarA] (4.47,0) rectangle (5.52,94.7);
      \node[PTval] at (3.805,82.4) {\PTrl{۸۲٫۴٪}};
      \node[PTval] at (4.995,94.7) {\PTrl{۹۴٫۷٪}};
      \node[PTcapt,anchor=north] at (4.4,-4) {\en{NDCG@3}};

      % ---- راهنما (از راست به چپ) --------------------------------------------
      \node[PTswatch,fill=PTaccentdark] at (9.30,116) {};
      \node[PTcapt,anchor=east]         at (9.08,116) {\PTrl{با بازرتبه‌بندی}};
      \node[PTswatch,fill=PTline,draw=PTmuted,line width=.4pt]
                                        at (6.10,116) {};
      \node[PTcapt,anchor=east]         at (5.88,116) {\PTrl{بدون بازرتبه‌بندی}};
    \end{tikzpicture}
  \end{center}

  \begin{itemize}
    \item \en{Precision@3} حدود \hl{۲۵٫۶ واحد درصد} بالا می‌رود؛ یعنی سه سند
          نهایی تقریباً همیشه مرتبط‌اند.
    \item رشد کمترِ \en{NDCG@3} نشان می‌دهد بازیابی پایه اسناد درست را
          \term{پیدا} می‌کرده، اما \term{مرتب} نمی‌کرده است.
  \end{itemize}
\end{frame}

%%%%%%%%%%%%%%%%%%%%%%%%%%%%%%%%%%%%%%%%%%%%%%%%%%%%%%%%%%%%%%%%%%%%%%%%%%%%%%%%
\begin{frame}{تحلیل خطا}
  \begin{PTcons}
    \begin{itemize}
      \item دو پرسش نادرست، هر دو از دستهٔ \term{عددی}: منبع درست بازیابی شد اما
            عدد در جدول صفحه بود و از متن برداشت نشد.
      \item پرسش‌هایی که پاسخشان در وب پراکنده است، به بیش از سه سند نیاز دارند.
    \end{itemize}
  \end{PTcons}

  \begin{itemize}
    \item تأخیر میانگین هر پرسش \hl{۴٫۲ ثانیه} است که بیشترش صرف بازیابی زندهٔ
          وب می‌شود، نه بازرتبه‌بندی.
    \item هزینهٔ هر پرسش نسبت به سپردن همهٔ اسناد به مدل، حدود \hl{یک‌سوم} است.
  \end{itemize}
\end{frame}
